\documentclass[10pt]{beamer}
\usetheme{metropolis}
\usepackage{graphicx}
\usepackage{listings}
\usepackage{booktabs}
\usepackage{xcolor}
\usepackage{hyperref}

\definecolor{codebg}{HTML}{F5F5F5}
\definecolor{codegreen}{HTML}{2E7D32}
\definecolor{codegray}{HTML}{757575}
\definecolor{codeblue}{HTML}{1565C0}

\lstdefinestyle{rstyle}{
  language=R, backgroundcolor=\color{codebg}, basicstyle=\ttfamily\scriptsize,
  keywordstyle=\color{codeblue}\bfseries, stringstyle=\color{codegreen},
  commentstyle=\color{codegray}\itshape, breaklines=true, frame=single,
  rulecolor=\color{codegray!30}, showstringspaces=false, tabsize=2,
  morekeywords={library,ggplot,aes,geom_point,geom_line,geom_col,geom_bar,
    geom_histogram,geom_boxplot,geom_smooth,geom_density,geom_jitter,
    scale_color_manual,scale_fill_brewer,scale_y_continuous,
    facet_wrap,coord_flip,theme_minimal,theme,element_text,element_blank,
    labs,stat_summary,position_dodge,position_stack,position_fill,position_jitter,margin}
}
\lstdefinestyle{pythonstyle}{
  language=Python, backgroundcolor=\color{codebg}, basicstyle=\ttfamily\scriptsize,
  keywordstyle=\color{codeblue}\bfseries, stringstyle=\color{codegreen},
  commentstyle=\color{codegray}\itshape, breaklines=true, frame=single,
  rulecolor=\color{codegray!30}, showstringspaces=false, tabsize=4,
  morekeywords={import,from,as,ggplot,aes,geom_point,geom_line,geom_col,
    geom_smooth,geom_histogram,geom_boxplot,geom_bar,geom_density,geom_jitter,
    scale_color_brewer,facet_wrap,theme_minimal,labs,coord_flip,
    position_dodge,position_stack,position_jitter,stat_summary}
}

\title{Wickham's Layered Grammar \& ggplot2}
\subtitle{Workshop 2 --- Module 2}
\author{Prof.Asoc.\ Endri Raco}
\institute{Data Visualization for Data Scientists \\ UNYT Tirana}
\date{2025/26}

\begin{document}

\begin{frame}
  \titlepage
\end{frame}

\begin{frame}{Module Roadmap}
  \begin{enumerate}
    \item The ggplot2 template: the pattern behind every plot
    \item \texttt{aes()}: mapping vs setting
    \item The geom catalogue and default statistics
    \item Position adjustments: dodge, stack, fill, jitter
    \item Aesthetic inheritance: global vs local
    \item Complete worked example in R and Python
  \end{enumerate}
  \begin{exampleblock}{Learning Outcome}
    You will master the ggplot2 syntax and construct any chart by composing layers with \texttt{+}.
  \end{exampleblock}
\end{frame}

\section{The ggplot2 Template}

\begin{frame}{The Pattern Behind Every Plot}
  \begin{center}
    \includegraphics[width=0.85\textwidth]{figures_w02m02/ggplot_template}
  \end{center}
\end{frame}

\begin{frame}{The \texttt{+} Operator Pipeline}
  \begin{center}
    \includegraphics[width=0.95\textwidth]{figures_w02m02/plus_operator}
  \end{center}
  \begin{alertblock}{Pro-Tip}
    The \texttt{+} must come at the \emph{end} of each line in R. In plotnine, place \texttt{+} at the \emph{start} inside parentheses.
  \end{alertblock}
\end{frame}

\section{Mapping vs Setting}

\begin{frame}{The \#1 ggplot2 Confusion}
  \begin{center}
    \includegraphics[width=0.92\textwidth]{figures_w02m02/mapping_vs_setting}
  \end{center}
\end{frame}

\begin{frame}[fragile]{Mapping vs Setting in Code}
  \begin{columns}[T]
    \begin{column}{0.48\textwidth}
      \textbf{MAPPING} (data-driven)
\begin{lstlisting}[style=rstyle]
# Colour varies with data
ggplot(df, aes(x = x, y = y,
    color = group)) +
  geom_point(size = 2)

# Size varies with data
ggplot(df, aes(x = x, y = y,
    size = population)) +
  geom_point(alpha = 0.6)
\end{lstlisting}
      {\small \texttt{aes()} creates a legend automatically.}
    \end{column}
    \begin{column}{0.48\textwidth}
      \textbf{SETTING} (fixed)
\begin{lstlisting}[style=rstyle]
# Colour fixed for all points
ggplot(df, aes(x = x, y = y)) +
  geom_point(color = "red",
             size = 3)

# Common mistake:
ggplot(df, aes(x = x, y = y,
    color = "red")) +
  geom_point()
# Maps string "red" as factor!
\end{lstlisting}
      {\small Fixed value inside \texttt{aes()} creates a useless single-level legend.}
    \end{column}
  \end{columns}
\end{frame}

\begin{frame}[fragile]{Mapping vs Setting in plotnine}
  \begin{columns}[T]
    \begin{column}{0.48\textwidth}
      \textbf{MAPPING}
\begin{lstlisting}[style=pythonstyle]
from plotnine import *

(ggplot(df, aes(x="x", y="y",
    color="group"))
 + geom_point(size=2))
\end{lstlisting}
    \end{column}
    \begin{column}{0.48\textwidth}
      \textbf{SETTING}
\begin{lstlisting}[style=pythonstyle]
from plotnine import *

(ggplot(df, aes(x="x", y="y"))
 + geom_point(color="red",
              size=3))
\end{lstlisting}
    \end{column}
  \end{columns}
  \vspace{4pt}
  \begin{exampleblock}{Data Insight}
    The rule is identical in R and Python: data-driven $\rightarrow$ inside \texttt{aes()}; fixed $\rightarrow$ outside \texttt{aes()}.
  \end{exampleblock}
\end{frame}

\section{The Geom Catalogue}

\begin{frame}{Eight Essential Geometries}
  \begin{center}
    \includegraphics[width=0.95\textwidth]{figures_w02m02/geom_catalogue}
  \end{center}
\end{frame}

\begin{frame}{Default Statistics per Geometry}
  \begin{center}
    \includegraphics[width=0.85\textwidth]{figures_w02m02/default_stats}
  \end{center}
  \begin{alertblock}{Pro-Tip}
    \texttt{geom\_bar()} counts rows (\texttt{stat\_count}); \texttt{geom\_col()} plots raw heights (\texttt{stat\_identity}). Confusing them is the \#2 ggplot2 error.
  \end{alertblock}
\end{frame}

\section{Position Adjustments}

\begin{frame}{Four Ways to Place Marks}
  \begin{center}
    \includegraphics[width=0.95\textwidth]{figures_w02m02/position_adjustments}
  \end{center}
\end{frame}

\begin{frame}[fragile]{Position Adjustments in R}
\begin{lstlisting}[style=rstyle]
# Stack (default for geom_bar with fill)
ggplot(df, aes(x = category, fill = type)) +
  geom_bar(position = "stack")

# Dodge (side by side)
ggplot(df, aes(x = category, fill = type)) +
  geom_bar(position = position_dodge(width = 0.8))

# Fill (proportional, sums to 100%)
ggplot(df, aes(x = category, fill = type)) +
  geom_bar(position = "fill") +
  scale_y_continuous(labels = scales::percent)

# Jitter (random noise to prevent overplotting)
ggplot(df, aes(x = category, y = rating)) +
  geom_jitter(width = 0.2, height = 0, alpha = 0.4, size = 1)
\end{lstlisting}
\end{frame}

\begin{frame}[fragile]{Position Adjustments in plotnine}
\begin{lstlisting}[style=pythonstyle]
from plotnine import *

# Stack
(ggplot(df, aes(x="category", fill="type")) + geom_bar(position="stack"))

# Dodge
(ggplot(df, aes(x="category", fill="type"))
 + geom_bar(position=position_dodge(width=0.8)))

# Fill
(ggplot(df, aes(x="category", fill="type")) + geom_bar(position="fill"))

# Jitter
(ggplot(df, aes(x="category", y="rating"))
 + geom_jitter(width=0.2, height=0, alpha=0.4, size=1))
\end{lstlisting}
\end{frame}

\section{The Statistics Layer}

\begin{frame}{Statistics Transform Data Before Rendering}
  \begin{center}
    \includegraphics[width=0.92\textwidth]{figures_w02m02/stat_layer}
  \end{center}
  \begin{exampleblock}{Data Insight}
    Every geom has a default stat, and every stat a default geom. You can override: \texttt{geom\_bar(stat = "identity")} or \texttt{stat\_smooth(geom = "point")}.
  \end{exampleblock}
\end{frame}

\begin{frame}[fragile]{stat\_summary: Compute on the Fly}
  \begin{columns}[T]
    \begin{column}{0.48\textwidth}
      \textbf{R}
\begin{lstlisting}[style=rstyle]
ggplot(apps, aes(x = Category,
                  y = Rating)) +
  stat_summary(
    fun = mean,
    fun.min = function(x)
      mean(x)-sd(x)/sqrt(length(x)),
    fun.max = function(x)
      mean(x)+sd(x)/sqrt(length(x)),
    geom = "pointrange",
    color = "#E53935") +
  coord_flip() +
  theme_minimal()
\end{lstlisting}
    \end{column}
    \begin{column}{0.48\textwidth}
      \textbf{Python (plotnine)}
\begin{lstlisting}[style=pythonstyle]
(ggplot(apps, aes(
    x="Category", y="Rating"))
 + stat_summary(
     fun_y=np.mean,
     fun_ymin=lambda x:
       np.mean(x)-np.std(x)/
       np.sqrt(len(x)),
     fun_ymax=lambda x:
       np.mean(x)+np.std(x)/
       np.sqrt(len(x)),
     geom="pointrange",
     color="#E53935")
 + coord_flip()
 + theme_minimal())
\end{lstlisting}
    \end{column}
  \end{columns}
\end{frame}

\section{Aesthetic Inheritance}

\begin{frame}{Global vs Local Aesthetics}
  \begin{center}
    \includegraphics[width=0.92\textwidth]{figures_w02m02/inheritance}
  \end{center}
\end{frame}

\begin{frame}[fragile]{Inheritance Rules}
  \begin{columns}[T]
    \begin{column}{0.48\textwidth}
      \textbf{Global} (in \texttt{ggplot()})
\begin{lstlisting}[style=rstyle]
ggplot(df, aes(x = displ,
    y = hwy, color = class)) +
  geom_point() +
  geom_smooth()
# Both inherit color = class
# -> one smooth line per class
\end{lstlisting}
    \end{column}
    \begin{column}{0.48\textwidth}
      \textbf{Local} (in \texttt{geom\_*()})
\begin{lstlisting}[style=rstyle]
ggplot(df, aes(x = displ,
    y = hwy)) +
  geom_point(aes(color = class)) +
  geom_smooth(color = "grey")
# Only points coloured by class
# Smooth = one overall grey line
\end{lstlisting}
    \end{column}
  \end{columns}
  \vspace{4pt}
  \begin{alertblock}{Pro-Tip}
    Place shared aesthetics in \texttt{ggplot(aes())} and layer-specific aesthetics in \texttt{geom\_*(aes())}. This controls which layers ``see'' a grouping variable.
  \end{alertblock}
\end{frame}

\section{Complete Worked Example}

\begin{frame}[fragile]{Complete Example: R}
  \begin{columns}[T]
    \begin{column}{0.52\textwidth}
\begin{lstlisting}[style=rstyle]
ggplot(mpg,
  aes(x = displ, y = hwy,
      color = class)) +
  geom_point(size = 2, alpha = 0.6) +
  geom_smooth(
    aes(group = 1),
    method = "lm", se = TRUE,
    color = "grey40") +
  scale_color_brewer(
    palette = "Set2") +
  labs(
    title = "Displacement vs MPG",
    x = "Engine (L)",
    y = "Highway MPG",
    color = "Class") +
  theme_minimal(base_size = 11)
\end{lstlisting}
    \end{column}
    \begin{column}{0.45\textwidth}
      \includegraphics[width=\textwidth]{figures_w02m02/complete_r}

      \vspace{4pt}
      {\small \texttt{aes(group=1)} overrides the global colour grouping in the smooth, producing one overall regression.}
    \end{column}
  \end{columns}
\end{frame}

\begin{frame}[fragile]{Complete Example: Python (plotnine)}
  \begin{columns}[T]
    \begin{column}{0.52\textwidth}
\begin{lstlisting}[style=pythonstyle]
from plotnine import *

(ggplot(mpg,
    aes(x="displ", y="hwy",
        color="class"))
 + geom_point(size=2, alpha=0.6)
 + geom_smooth(
     aes(group=1),
     method="lm", se=True,
     color="grey", size=1)
 + scale_color_brewer(
     type="qual", palette="Set2")
 + labs(
     title="Displacement vs MPG",
     x="Engine (L)",
     y="Highway MPG",
     color="Class")
 + theme_minimal())
\end{lstlisting}
    \end{column}
    \begin{column}{0.45\textwidth}
      \includegraphics[width=\textwidth]{figures_w02m02/complete_py}

      \vspace{4pt}
      {\small Near-identical code. Differences: strings for columns, \texttt{+} at line start, \texttt{size} vs \texttt{linewidth}.}
    \end{column}
  \end{columns}
\end{frame}

\section{Synthesis}

\begin{frame}{Key Takeaways}
  \begin{enumerate}
    \item Every ggplot2 chart follows the same \textbf{template}: \texttt{ggplot + geom + scale + facet + labs + theme}.
    \item \textbf{Mapping} (inside \texttt{aes}) is data-driven; \textbf{setting} (outside \texttt{aes}) is fixed.
    \item Each geom has a \textbf{default stat}: \texttt{geom\_bar} counts; \texttt{geom\_col} plots identity.
    \item \textbf{Position adjustments} (dodge, stack, fill, jitter) control how marks overlap.
    \item \textbf{Global} aesthetics are inherited by all layers; \textbf{local} apply to one layer only.
    \item \textbf{plotnine} syntax is near-identical to ggplot2.
  \end{enumerate}
\end{frame}

\begin{frame}{Essential Readings}
  \begin{itemize}
    \item Wickham, H. (2010). ``A Layered Grammar of Graphics.'' \emph{JCGS}, 19(1), 3--28.
    \item Wickham, H. (2016). \emph{ggplot2}, 2nd ed., Chapters 2--3.
    \item Wickham, H. \& Grolemund, G. (2023). \emph{R for Data Science}, 2nd ed., Chapters 9--12.
    \item plotnine documentation: \texttt{https://plotnine.org}
  \end{itemize}
  \vspace{6pt}
  \textbf{Next Module}: ggplot2 Deep Dive --- Aesthetics \& Geometries (W02-M03)
\end{frame}

\end{document}
